\documentclass[10pt,a4paper]{article}
\usepackage[margin=2cm]{geometry}
\usepackage{needspace}
\usepackage{fontspec,booktabs,longtable,array,xcolor,hyperref}
\setmainfont{TeX Gyre Pagella}
\setsansfont{TeX Gyre Heros}
\definecolor{reportblue}{HTML}{244D65}
\hypersetup{colorlinks=true,urlcolor=reportblue,linkcolor=reportblue}
\setlength{\parindent}{0pt}
\setlength{\parskip}{5pt}
\renewcommand{\arraystretch}{1.18}
\setlength{\tabcolsep}{5pt}
\emergencystretch=2em
\urlstyle{same}
\pagestyle{plain}
\title{\sffamily Jev on SocialIQA: Social Commonsense Decisions}
\author{}
\date{September 30, 2026}
\begin{document}
\begin{center}{\sffamily\LARGE Jev on SocialIQA: Social Commonsense Decisions\par}\vspace{8pt}{\small September 30, 2026}\end{center}\vspace{6pt}
\section*{Abstract}
Jev answered 1,792 of 2,224 English SocialIQA v1.4 test questions correctly, achieving 80.58\% accuracy with a duplicate-question cluster 95\% interval of 78.81\%–82.21\%. Accuracy across the nine official source groups ranges from 78.11\% to 84.66\%. These results describe selection judgments about everyday motives, emotions and consequences.
\section{Dataset}
The material contains all 2,224 rows of the authors’ v1.4 test release. Each item contains an English context, a question and three options. Each is answered independently and scored against the released reference label. Model snapshot: typesafe/jev-1.13-20260917.

The nine groups use the official promptDim field, which records the ATOMIC question-generation source. Rewritten questions can depart from that source, so these groups are not treated as nine independent psychological abilities. The evaluation concerns a long-public English multiple-choice dataset.

\section{Results}
\Needspace{8\baselineskip}
\subsection{Overall results}
The dataset labels are the scoring criterion. Every formal item has a valid response. Uniform random three-option selection has an expected accuracy of 33.33\%.

{\small
\begin{longtable}{@{}>{\raggedright\arraybackslash}p{\dimexpr 0.4000\textwidth-2\tabcolsep\relax}>{\raggedright\arraybackslash}p{\dimexpr 0.6000\textwidth-2\tabcolsep\relax}@{}}
\toprule
\textbf{Measure} & \textbf{Result}\\
\midrule
\endfirsthead
\toprule
\textbf{Measure} & \textbf{Result}\\
\midrule
\endhead
Formal items & 2,224\\
Correct & 1,792\\
Accuracy & 80.58\%\\
Question-cluster 95\% interval & 78.81\%–82.21\%\\
Distinct complete questions & 2203\\
Distinct-question accuracy & 80.57\%\\
\bottomrule
\end{longtable}
}
\newpage
\subsection{Official source groups}
{\small
\begin{longtable}{@{}>{\raggedright\arraybackslash}p{\dimexpr 0.3200\textwidth-2\tabcolsep\relax}>{\raggedright\arraybackslash}p{\dimexpr 0.2267\textwidth-2\tabcolsep\relax}>{\raggedright\arraybackslash}p{\dimexpr 0.2267\textwidth-2\tabcolsep\relax}>{\raggedright\arraybackslash}p{\dimexpr 0.2267\textwidth-2\tabcolsep\relax}@{}}
\toprule
\textbf{Group} & \textbf{Items} & \textbf{Correct} & \textbf{Accuracy}\\
\midrule
\endfirsthead
\toprule
\textbf{Group} & \textbf{Items} & \textbf{Correct} & \textbf{Accuracy}\\
\midrule
\endhead
oEffect (Effect on others) & 159 & 127 & 79.87\%\\
oReact (Others’ reactions) & 223 & 179 & 80.27\%\\
oWant (Others’ wants) & 252 & 201 & 79.76\%\\
xAttr (Actor attributes) & 297 & 232 & 78.11\%\\
xEffect (Effect on actor) & 103 & 82 & 79.61\%\\
xIntent (Actor intent) & 297 & 239 & 80.47\%\\
xNeed (Actor prerequisites) & 277 & 225 & 81.23\%\\
xReact (Actor reactions) & 277 & 220 & 79.42\%\\
xWant (Actor wants) & 339 & 287 & 84.66\%\\
\bottomrule
\end{longtable}
}
The xWant group has the highest accuracy and xAttr the lowest. Group sizes differ; this ordering describes the dataset rather than a stable hierarchy of abilities.

\Needspace{8\baselineskip}
\subsection{Call costs}
{\small
\begin{longtable}{@{}>{\raggedright\arraybackslash}p{\dimexpr 0.2400\textwidth-2\tabcolsep\relax}>{\raggedright\arraybackslash}p{\dimexpr 0.3800\textwidth-2\tabcolsep\relax}>{\raggedright\arraybackslash}p{\dimexpr 0.3800\textwidth-2\tabcolsep\relax}@{}}
\toprule
\textbf{Input tokens} & \textbf{Output tokens} & \textbf{Reported cost (USD)}\\
\midrule
\endfirsthead
\toprule
\textbf{Input tokens} & \textbf{Output tokens} & \textbf{Reported cost (USD)}\\
\midrule
\endhead
851,327 & 84,512 & 0.035755734\\
\bottomrule
\end{longtable}
}
Reported formal-evaluation cost is USD 0.035755734. Twelve separate interface-debugging items bring the reported total to USD 0.035949816, with 855,948 input tokens and 84,968 output tokens. One technical failure has unknown billing; the historical ledger retains USD 0.01 as a reservation, not a settled charge.

\Needspace{8\baselineskip}
\section{Conclusion}
Jev achieves 80.58\% accuracy on this English social-commonsense dataset, selecting the reference answer for most questions about motives, emotions and consequences. It misses 432 reference answers. The result provides a baseline for social-situation understanding under the tested selection format; it does not measure emotional experience or willingness to help.
\Needspace{5\baselineskip}
\section*{References}
\begin{enumerate}
\item Sap et al. (2019). Social IQa: Commonsense Reasoning about Social Interactions. EMNLP-IJCNLP. \url{https://aclanthology.org/D19-1454/}
\item SocialIQA v1.4: authors’ dataset release, CC BY 4.0. \url{https://maartensap.com/social-iqa/}
\end{enumerate}
\end{document}
